\section{Introduction}
\label{sec:intro}

Contracts are amended by later agreements, regulations by later notifications, and technical standards by later
standards. An answer drawn from such a collection can quote a provision correctly and still be wrong, because the
document it quotes was replaced before the date the question is about. Systems that use language models
for this can add a second model, an LLM judge, to check the answers before they reach a user. The judge is
often given help: an amendment record listing which document replaced which, and from when.

Our starting point is a practical worry about that arrangement. Amendment records are seldom complete. Consolidated
versions lag behind new notifications; pipelines that extract links from documents miss some; registers cover only the
instruments they were built from. Wherever a record is silent about a replacement, it effectively reports a replaced
document as current. If the judge reads the documents in front of it, such gaps cost little. If it relies on the record,
it will pass exactly the answers the record cannot catch, and it adds little beyond what a lookup in the record would do.

Existing evaluations of LLM judges do not tell a builder which of these to expect. Benchmarks give the judge correct
context or none \citep{tan2025judgebench,xu2025contextualjudgebench}, and most report whether the judge flagged an
answer, not why. A judge that flags every older citation would score well on detection, and so would one that flags
any older document that sits next to a newer one on the same subject. Neither is checking currency.

We therefore built a test aimed at one error, citing a replaced document, and at the question of what the judge relies
on when the record and the documents disagree. The test starts from any table of replacements. For each replacement it
writes an answer citing the replaced document (a fault) and three answers that are correct: one citing the successor,
one citing the replaced document at a date before the replacement existed (a date control), and one citing an older
document that a newer one updated without replacing (a decoy; Figure~\ref{fig:design} shows one link). It then removes half the links from the record, presents
the damaged record in several ways that leave its content unchanged, and counts a fault as caught only when the judge
says the cited document was replaced.

We applied the test to \ENLinks{} replacements recorded by the RFC Editor of the Internet Engineering Task Force (IETF), with four judges (Gemini~3.1~Pro,
GPT-5.5, Claude Opus~4.5 and the open-weight Gemma~4~26B). The hypotheses, tests and the way we would read each outcome
were committed to the project's private development repository (the release lists the commit times) after a smaller pilot (mostly on 24 of the links) and before these runs.

With a correct record line, the judges caught every replaced citation or nearly so, which a lookup in the record
also does. On the faults whose link we removed, every judge caught fewer with the record present than from the texts
alone, on each of two random choices of which links to remove, and most of the missed answers were passed without
any finding. The drop also occurred when the line said only that no replacement was recorded, a statement that is
true. Presenting the record as an unverified note, or warning that it was incomplete, returned Gemini, GPT-5.5 and
Claude most of the way to their text-only behaviour, false alarms included; the change is significant after correction
for Gemini and GPT-5.5 but not for Claude, which had fewer catches to recover, and Gemma recovered at most \EGemmaArmCaveatGained{} of the \EGemmaHTwoaLost{}
it had lost (\S\ref{sec:res-gaps}). Offered a reason code for a record that conflicts with the evidence, in an
exploratory context added afterwards, Gemini named the conflict on \EGeminiConflictCodeGap{} of the 24 gap faults and
the other judges rarely or never did, so the silence is partly a matter of reply format for one judge only. Given both texts and no record, the judges caught many more replacements than they
flagged updates, but every update they flagged had a near-identical title, and an unplanned analysis suggests that
they treat the resemblance between two documents as evidence of replacement.
In the pilot, date controls exposed a 32B open model that flagged almost every old citation whatever its date; the
judges here passed nearly all of them.

The paper's contribution is the test, released as a small harness, together with this evidence about four judges on
one corpus. We do not claim that LLM judges in general defer to records, and we cannot say why the judges we tested
did. A second setting, Indian rooftop-solar regulation with an amendment record extracted by an AI pipeline, is
reported as a descriptive case study in the appendix: its answer key did not pass our own quality checks, and nothing
in the main results rests on it.
